\section{Geometric exclusions for a nonvanishing counterexample}
\label{proj:sec:exclusions}

Throughout this section we assume the lower-dimensional good-model
hypothesis of Assumption~\ref{proj:mmp:lower-models}. We suppose, towards a
contradiction, that \(X\) is smooth projective of dimension \(n>0\) and
\begin{equation}\label{proj:ex:counterexample}
 K_X\ \text{is pseudo-effective},
 \qquad \kappa(X,K_X)=-\infty.
\end{equation}
These properties persist on smooth projective birational models. All
varieties in this section are complex. Numerical dimension for a
pseudo-effective divisor means Nakayama's numerical dimension
\(\kappa_\sigma\); for a nef divisor it agrees with the intersection
definition used in Theorem~\ref{proj:thm:signed}.

\subsection{The Albanese reduction and algebraic webs}

\begin{lemma}\label{proj:ex:irregularity}
Every smooth projective birational model of \(X\) has irregularity zero.
Consequently, on such a model, or on a projective \(\Q\)-factorial
terminal birational model, numerical equivalence of rational divisors
implies their \(\Q\)-linear equivalence.
\end{lemma}
\begin{proof}
If \(q(X)>0\), resolve the Stein factorization of the Albanese map to
obtain a morphism \(f:W\to B\) with connected fibers, with \(W,B\)
smooth projective and \(\dim B>0\). The map from \(B\) to a subvariety
of \(\Alb(X)\) is generically finite. Wedges of invariant one-forms
on the abelian variety therefore give a nonzero section of \(K_B\):
at the generic point, choose \(\dim B\) independent pulled-back
one-forms. Thus \(\kappa(B,K_B)\geq0\).

For a very general smooth fiber \(F\), the restriction of the
pseudo-effective class \(K_W\) is pseudo-effective. For example, restrict
a positive current representing it to almost every fiber, or use
effective approximations with arbitrarily small ample error.
Adjunction identifies this restriction with \(K_F\). If \(F\) has
positive dimension, the induction hypothesis gives
\(\kappa(F,K_F)\geq0\); for a point the same assertion holds by
convention. Applying Assumption~\ref{proj:asm:iitaka} with both boundaries
empty gives \(\kappa(W,K_W)\geq0\), contradicting
\eqref{proj:ex:counterexample}.

This is the only use of Assumption~\ref{proj:asm:iitaka} in the proof.
Irregularity is a smooth birational invariant. Finally, when
\(\Pic^0=0\), a numerically trivial line bundle is torsion, since the
group of numerically trivial line bundles modulo \(\Pic^0\) is
finite. Clear denominators for rational divisors. On a terminal
\(\Q\)-factorial model, pull back to a smooth resolution and then
descend the rational linear equivalence.
\end{proof}

We use covering families of proper subvarieties through very general
points. They can be parameterized in countably many algebraic
families, using Hilbert schemes followed by resolution and
stratification. After shrinking a parameter space, the domains form
a smooth projective family with integral fibers and have a dominant
evaluation map. If the fiber maps are generically finite onto their
images, general ample cuts of the parameter space make the total
evaluation generically finite and still dominant. We always resolve
and compactify this evaluation when applying ramification.
Invariance of plurigenera for smooth projective families
\cite[Theorem~1]{Proj-Paun2007} permits the same construction for the
Iitaka fibers of general members: choose a divisible pluricanonical
system on the general member, spread it by base change, and resolve
its relative rational map.

\begin{lemma}[Algebraic web quotient]\label{proj:ex:web}
Let a smooth projective variety be dominated generically finitely by
the total space of a family with smooth integral projective general
fibers. On a dense regular open, form the distribution generated by
the tangent spaces of the fiber images, taking spans, saturation,
and Lie brackets. Its general leaves are dense opens of algebraic
subvarieties. Their closures are the general fibers of a rational
map, which has connected general fiber after Stein factorization.
\end{lemma}
\begin{proof}
Shrink to an open \(U\) where the evaluation has finitely many etale
sheets, the parameter map is smooth, and the generated involutive
distribution \(\mathcal F\) has constant rank. Its construction is
algebraic: spans and brackets stabilize at the generic point after
finitely many operations, and descend from the etale sheets. The
nonempty open parts of the integral parameter fibers are connected.
We use only chains of fiber-image steps lying in \(U\).

For \(x\in U\), endpoints of chains of at most \(k\) steps form a
constructible set, by the algebraic fiber-product construction and
Chevalley's theorem. Every endpoint lies in the analytic leaf
through \(x\). A locally closed smooth variety contained in an
analytic leaf has dimension at most \(\rk\mathcal F\). Indeed, in a
Frobenius chart the leaf meets at most countably many plaques: use
finite plaque chains in a countable foliation atlas. The transverse
coordinate map on any connected smooth piece then has countable
image and is constant.

There is therefore a maximal dimension among the closures of all
such endpoint loci; it is attained for some finite \(k\). The loci
are nested because we allow at most \(k\) steps, including the
identity step. Choose an irreducible component \(V\) of maximal
dimension and a dense open \(O\subset V\) of reachable points.
For each etale sheet, restrict its fiber-equivalence relation to
first coordinate in \(O\), and take the component through the
diagonal section. Its second-image closure contains \(O\), hence
\(V\), while its second image consists of endpoints reachable
in at most \(k+1\) steps. Maximality forces that closure to
equal \(V\). At a general
diagonal point, smoothness of the parameter map identifies its
vertical tangent with the corresponding web tangent. Every web
tangent is consequently tangent to \(V\), and so is every bracket. Hence
\(\dim V\geq\rk\mathcal F\), while the reverse inequality was proved
above. A plaque is thus open in \(V\). The equations of \(V\),
pulled back to the connected immersed leaf, vanish on a nonempty
open and hence on the entire leaf. Its closure is exactly \(V\).

These rank-dimensional invariant subvarieties have countably many
Hilbert or Chow parameter spaces. Tangency on the regular locus is
an algebraic condition after stratification, so one such family
dominates. An invariant irreducible variety meeting \(U\) contains
the leaf through any of its points in \(U\): the tangent vector
fields preserve its reduced ideal, first on its smooth locus and
then everywhere by closure. A rank-dimensional leaf closure is
therefore unique through a general point. Assigning that closure
gives the desired rational map. Resolve its graph and take the
Stein factorization.
\end{proof}

For positive closed \((1,1)\)-currents we shall use this quotient in
the following way. If a current vanishes on almost every parameter
fiber away from a fixed removed divisor, it annihilates the
corresponding fiber tangents on a dense open. In submersion
coordinates, choose locally integrable plurisubharmonic potentials.
Fubini's theorem makes the pure fiber Hessian zero as a
distribution; positivity then makes the mixed coefficients in
fiber directions zero as well. A generically etale evaluation
transfers this assertion to the target. Annihilation passes to
brackets by the identities
\[
 \mathcal L_v T=d(\iota_vT)+\iota_v(dT),\qquad
 \iota_{[v,w]}T=\mathcal L_v(\iota_wT)-\iota_w(\mathcal L_vT).
\]
Pullbacks by dominant morphisms and restrictions to almost every
smooth parameter fiber are defined by the same local potentials.

\subsection{The general-type exclusion}

\begin{proposition}\label{proj:prop:general-type}
Every positive-dimensional proper subvariety through a very general
point of \(X\) is of general type on resolution. The assertion holds
also on every projective birational model.
\end{proposition}
\begin{proof}
Otherwise take a covering family of nongeneral-type subvarieties,
with smooth domains \(V\), and slice its parameters as above.
Restricting the ramification formula of its generically finite
total evaluation shows that \(K_V\) dominates the restriction of
the pulled-back \(K_X\) by an effective divisor. Thus \(K_V\) is
pseudo-effective. The induction hypothesis gives a good minimal
model of \(V\), so \(\kappa(V)\geq0\). Since \(V\) is not of general
type, its general Iitaka fibers have positive dimension and
Kodaira dimension zero. Spreading these fibers gives another
covering family with smooth domains \(G\), of dimension less than
\(n\), and \(\kappa(G)=0\). Each \(G\) has a good minimal model and
\(\kappa_\sigma(K_G)=0\).

Fix a positive current \(T\) representing \(K_X\). After slicing and
resolving the total evaluation of the \(G\)'s, its restriction to
almost every parameter fiber, plus the effective restricted
ramification divisor, is a positive current in \(K_G\).
Every positive current in that class is supported on the fixed
canonical divisor of \(G\). To see this, take a common resolution
with its good minimal model. The latter has torsion canonical
divisor, so the canonical divisor upstairs is effective exceptional.
Intersection with a pullback of an ample class to the power
\(\dim G-1\) forces any such positive current to vanish off the
exceptional locus. The support theorem makes it divisorial, and
independence of exceptional divisor classes, by negativity, fixes
its coefficients. Pushing down proves the assertion on \(G\).
The excluded support is algebraic in the family: it is the fixed
divisor of a sufficiently divisible relative pluricanonical
system after shrinking the parameter space.

It follows that \(T\) annihilates the web distribution of the
\(G\)'s on a dense open. If that distribution has full rank, \(T\)
is supported on a proper algebraic set. The support theorem for
positive closed \((1,1)\)-currents expresses it as a finite
nonnegative real combination of prime divisors
\cite{Proj-Siu1974,Proj-Demailly2012}. Its rational cohomology class then has
a nonnegative rational representative: solve the finite rational
linear system for the coefficients on the same face of the
nonnegative orthant. Lemma~\ref{proj:ex:irregularity} converts numerical
effectivity to \(\Q\)-linear effectivity, a contradiction.

Otherwise Lemma~\ref{proj:ex:web} gives a rational quotient with general
smooth fiber \(F\) on a smooth resolved graph, where
\(0<\dim F<n\). Its canonical divisor is pseudo-effective.
The moving \(G\)'s still generate its tangent space at general
points. Indeed the quotient is constant on every general web
member, hence induces a rational map on their parameter space;
restricting to a general quotient fiber preserves dominance of
evaluation. A pluricanonical rational map of \(F\) is constant
along each such \(G\). Slice the parameters inside \(F\) to make
the total evaluation generically finite: ramification injects the
restricted pluricanonical sections into sections of a multiple of
\(K_G\), and these span a space of dimension at most one.
Consequently their ratios are constant on \(G\). Their
differentials vanish on the web and its brackets, so the
pluricanonical map is constant on \(F\). Lower-dimensional
nonvanishing therefore gives \(\kappa(F)=0\), and its good minimal
model gives \(\kappa_\sigma(K_F)=0\).

Now apply the numerical-zero-fiber reduction of
Gongyo--Lehmann \cite[Theorem~1.3 and Corollary~4.5]{Proj-GongyoLehmann2013}.
For a projective \(\Q\)-factorial klt rational pair with a
connected-fiber morphism and numerical dimension zero on the
general fiber, that theorem produces a klt pair on a smooth
birational base whose good-minimal-model existence is equivalent.
Here the total space is the smooth graph, the boundary is zero,
and the base has dimension less than \(n\). The required numerical
dimension is \(\kappa_\sigma\), following the convention in
\cite[Section~3 and Theorem~3.2]{Proj-Fujino2019Numerical}. It is
zero by the good model of \(F\), so the numerical-dimension
hypothesis holds. The induction hypothesis supplies the base good
model, hence nonvanishing upstairs, contradicting
\eqref{proj:ex:counterexample}.
\end{proof}

\begin{corollary}\label{proj:ex:supporting}
Let \(Y\) be a projective \(\Q\)-factorial terminal birational model
of \(X\). If a rational divisor \(P\) has \(\kappa(Y,P)\geq1\),
then \(K_Y+cP\) is big for every rational \(c>0\).
Consequently, if \(\alpha\neq0\) is a supporting functional of the
pseudo-effective cone with \(\alpha(K_Y)=0\), then
\(\alpha(P)>0\).
\end{corollary}
\begin{proof}
Resolve a moving subsystem of a multiple of \(P\). If its map is
generically finite, \(P\) is big and the assertion follows from
pseudoeffectivity of \(K_Y\). Otherwise its general fiber is a
proper positive-dimensional subvariety through very general
points and is of general type by Proposition~\ref{proj:prop:general-type}.
The canonical divisor of the smooth resolution is relatively big,
so adding a sufficiently large ample pullback from the image makes
it big. Since that pullback is bounded by a multiple of the
resolved moving subsystem, pushforward gives \(K_Y+aP\) big for
some \(a>0\). Convexity with the pseudo-effective \(K_Y\), and then
addition of the pseudo-effective \(P\), gives the assertion for
every \(c>0\). A nonzero nonnegative functional on a closed convex
cone is strictly positive on its interior. Applying it to
\(K_Y+cP\) proves the last assertion.
\end{proof}

\subsection{Valuations and positive currents}

For a positive closed \((1,1)\)-current \(T\) on a smooth projective
variety \(V\), let \(\nu_E(T)\) denote the generic Lelong number of
its pullback at a divisorial valuation \(E\). The normalization is
such that a reduced smooth divisor has generic number one. We
write \(A_V(E)=a(E;V,0)\) for log discrepancy.

\begin{lemma}\label{proj:ex:valuation-acc}
For a fixed \(T\) and a fixed \(M\), the set
\[
 \{\nu_E(T): A_V(E)\leq M\}
\]
satisfies the ascending chain condition.
\end{lemma}
\begin{proof}
We induct on the integer part of \(M\); discrepancies over a
smooth variety are positive integers. Suppose that a strictly
increasing sequence exists, and discard initial terms so its
members exceed a fixed positive number \(c\).
Skoda integrability and change of variables give
\begin{equation}\label{proj:ex:skoda-valuation}
 \nu_E(T)\leq A_V(E)\nu_x(T)
\end{equation}
at a general point \(x\) of the center. Indeed every exponent
smaller than \(1/\nu_x(T)\) is locally integrable, whereas
integrability after pullback imposes the corresponding
discrepancy bound. Hence all centers lie in the fixed proper Siu
locus \(\{\nu_x(T)\geq c/M\}\), which is algebraic by Siu
analyticity and projectivity \cite{Proj-Siu1974}.

If, after passage to a subsequence, the centers lie in a fixed
codimension-at-least-two subvariety, principalize its ideal.
All lifted centers lie in the exceptional locus, where the
relative canonical divisor has positive integral order.
The discrepancy bound for the new smooth ambient space has
therefore decreased by at least one. Pull back \(T\) and apply
induction.

Otherwise the centers lie in a fixed prime divisor \(D\) of the
Siu locus. Write \(T=c_D[D]+T'\) with \(T'\) positive and with
zero generic Lelong number along \(D\). The integers
\(\ord_E(D)\) are bounded: the fixed effective divisor \(D\) has
positive log canonical threshold, and
\(\operatorname{lct}(D)\ord_E(D)\leq A_V(E)\).
Pass to a constant order. The residual numbers \(\nu_E(T')\)
are still strictly increasing, and after discarding a term have
a positive lower bound. Equation~\eqref{proj:ex:skoda-valuation}
places their centers in a fixed Siu locus of \(T'\), which does
not contain \(D\). Their intersection with \(D\) has codimension
at least two, so the preceding case applies.
\end{proof}

We also record the multiplier-ideal approximation used below
\cite[Theorems~5.11, 6.27, and 14.2]{Proj-Demailly2012}. On a smooth projective variety, if \(\{T\}\)
is a real algebraic class, one can choose integral line bundles
\(L_k\) such that \(c_1(L_k)-k\{T\}\) stays in a bounded,
uniformly sufficiently ample set and
\(L_k\otimes\mathcal J(kT)\) is globally generated. Round the
coefficients of \(k\{T\}\) in a fixed integral basis and add a
fixed sufficiently positive integral class. Nadel vanishing and
Castelnuovo--Mumford regularity give generation. For every
divisorial valuation,
\begin{equation}\label{proj:ex:multiplier-order}
 \ord_E\mathcal J(kT)\leq k\nu_E(T).
\end{equation}
The local Bergman weight is bounded below by the original weight
up to a constant; this inequality persists after pullback and
gives \eqref{proj:ex:multiplier-order}. These are statements for a
fixed current, not uniform assertions over all currents.

\begin{proposition}\label{proj:prop:no-pullback}
Let \(Y\) be a fixed projective \(\Q\)-factorial terminal
birational model of \(X\), with \(K_Y\) nonbig. A positive current
in the pullback class of \(K_Y\) on a smooth resolution cannot,
on a dense regular Zariski open, annihilate the tangent spaces
of the general fibers of a rational fibration of positive
relative dimension.
\end{proposition}
\begin{proof}
Suppose otherwise, and choose a connected-fiber quotient of
minimal base dimension \(b\). If \(b=0\), the current vanishes on
a dense open; the divisorial-current argument in
Proposition~\ref{proj:prop:general-type} contradicts
\eqref{proj:ex:counterexample}. Thus \(0<b<n\).
Resolve the graph and the space carrying the current:
\[
 \begin{tikzcd}
 W \arrow[r,"p"] \arrow[d,"h"'] & Y \\
 B &
 \end{tikzcd}
\]
Here \(W,B\) are smooth projective and \(h\) has connected
general fiber. Flatten the main component over a modification of
the base, resolve that base, normalize the main transform, and
then resolve upstairs. Every vertical prime on the flat
transform maps to a divisor of the base: a prime over base
codimension \(c\) would have generic fiber dimension at least
\(n-b+c-1\), whereas all fibers have dimension at most \(n-b\).
Consequently every vertical divisor
on \(W\) over base codimension at least two is exceptional over
\(Y\). This property persists under later resolutions and base
modifications: a nonexceptional prime already maps onto a base
prime and its generic image is unchanged.

For divisors and numerical classes put
\[
 T_B=p_*h^*.
\]
This strict pullback is well defined on numerical classes because
\(Y\) is \(\Q\)-factorial: numerical classes upstairs decompose
into pulled-back classes from \(Y\) and exceptional classes.
It preserves pseudoeffectivity. For a higher base model
\(r:B'\to B\), its analogue satisfies
\(T_{B'}r^*=T_B\), by computing on a common graph.
It kills divisors exceptional over \(B\): otherwise a prime
in such a pullback that dominates a divisor of \(Y\) would
have center of codimension at least two on \(B\), contrary
to the property just established.

\smallskip
\noindent\emph{Descent of the current.}
Let \(T\) also denote the pulled-back current on \(W\).
Shrink to a smooth open \(B^0\) so that \(h\) is smooth
proper, the removed set contains no vertical divisor
over \(B^0\), and the remaining open in every fiber is
nonempty and connected. Away from that fixed removed
algebraic set upstairs, \(T\) descends to a positive current
\(S^0\).
Indeed in submersion coordinates the horizontal coefficient
distributions are independent of fiber variables by closedness,
and the open parts of general fibers are connected. They
therefore glue to \(S^0\) on \(B^0\). The difference
\(T-h^*S^0\) is closed of order zero and supported on the
removed algebraic set. The support theorem expresses it as a
signed divisor sum; a closed order-zero \((1,1)\)-current
supported in codimension at least two is zero. Its horizontal coefficients are
nonnegative, since a pullback has zero generic Lelong number
along a horizontal divisor.

Subtract those finitely many global Siu components from
\(T\), obtaining a positive current \(T'\) that equals
\(h^*S^0\) on all of \(h^{-1}(B^0)\). Choose a K\"ahler form
\(\omega\) representing an ample algebraic class and put
\(d=n-b\). Fiber integration gives
\(h_*(\omega^d)=c>0\), constant on \(B^0\), by closedness.
The projection formula therefore gives
\[
 c^{-1}h_*(T'\wedge\omega^d)|_{B^0}=S^0.
\]
The left side defines a global positive closed extension.
Its class is real algebraic, since \(\{T'\}\) is the
algebraic class \(p^*K_Y\) minus real divisor classes and
algebraic intersection and pushforward preserve such
classes. Remove its divisorial parts along
\(B\setminus B^0\), and call the result \(S\).
On a dense open \(T=h^*S\). Their difference is again closed
of order zero, so the same support theorem leaves only a
signed divisor sum. It has nonnegative coefficients at every
prime dominating a base prime. For the latter assertion, the
generic local map is a transverse power map times a
submersion. A current with zero generic Lelong number along
the base prime has zero generic number at such an upstairs
prime: local target balls of a radius equal to a fixed power
of the source radius give the usual Lelong comparison.
The only possible negative coefficients lie over base
codimension at least two and are exceptional over \(Y\).
It follows that
\begin{equation}\label{proj:ex:descent-domination}
 K_Y-T_B\{S\}\quad\text{is pseudo-effective}.
\end{equation}
Horizontal parts subtracted above merely contribute effective
divisors to this difference.

\smallskip
\noindent\emph{A negative canonical direction on the base.}
Choose a nonzero supporting functional \(\alpha\) of the
pseudo-effective cone at \(K_Y\), and put
\(\eta=T_B^*\alpha\). The corresponding functionals \(\eta'\)
on higher smooth base models are movable classes by
pseudo-effective/movable duality \cite[Theorem~0.2]{BDPP}.
They annihilate base-exceptional divisors. Equation
\eqref{proj:ex:descent-domination} gives \(\eta\cdot\{S\}=0\).

Only finitely many prime divisors on \(B\) have positive
divisorial Lelong coefficient for \(S\). Otherwise their
nonzero strict pullbacks are effective divisors on \(Y\)
killed by \(\alpha\), with pairwise disjoint prime supports.
The supports are disjoint because a nonexceptional prime of
\(Y\) cannot map into the intersection of two base primes.
Only finitely many base primes have zero strict pullback,
since their total pullbacks are supported in the finite
exceptional locus of \(p\). In the finite-dimensional rational
space of divisor classes, infinitely many remaining strict
pullbacks have a rational relation. Its positive and negative
sides are nonzero effective divisors with disjoint support;
Lemma~\ref{proj:ex:irregularity} makes the relation
\(\Q\)-linear. They define a moving pencil killed by
\(\alpha\), contradicting Corollary~\ref{proj:ex:supporting}.

We claim that
\begin{equation}\label{proj:ex:negative-base}
 K_B\cdot\eta<0.
\end{equation}
The smooth general fiber \(F\) of \(h\) is of general type by
Proposition~\ref{proj:prop:general-type}. The relative-positivity
theorem of Kov\'acs--Patakfalvi
\cite[Theorem~9.9]{Proj-KovacsPatakfalvi2017} applies to a log
canonical fiber space with smooth base, SNC total pair, and
log-general-type geometric generic fiber; for a rational base
divisor \(M\) with \(\kappa(M)\geq0\), it gives relative
subadditivity with the term \(h^*M\). Here both spaces are
smooth projective, the boundary is zero, and we take \(M=0\).
Invariance of plurigenera on the smooth locus identifies the
generic-fiber Kodaira dimension with that of \(F\). Thus
\begin{equation}\label{proj:ex:relative-positive}
 \kappa(W,K_W-h^*K_B)\geq \kappa(F,K_F)=n-b>0.
\end{equation}
This is an established general-type-fiber theorem, not another
use or strengthening of Assumption~\ref{proj:asm:iitaka}.
Choose compatible canonical divisors and set
\(P=K_Y-T_B(K_B)=p_*(K_W-h^*K_B)\).
For every sufficiently divisible \(m\), pushing forward an
effective divisor
\(\operatorname{div}_W(\varphi)+m(K_W-h^*K_B)\)
gives
\(\operatorname{div}_Y(\varphi)+mP\geq0\).
The resulting injection of sections preserves their ratios,
so \(\kappa(Y,P)\geq1\). Corollary~\ref{proj:ex:supporting} now gives
\(0<\alpha(P)=-K_B\cdot\eta\), proving
\eqref{proj:ex:negative-base}.

\smallskip
\noindent\emph{Rational curves and an exact zero gap.}
Apply the multiplier approximation to \(S\) on \(B\), and
principalize \(\mathcal J(kS)\) by \(r_k:B_k\to B\).
If \(F_k\) is its divisor, the class
\[
 P_k=\frac{r_k^*L_k-F_k}{k}
\]
is nef. Put \(\eta_k=T_{B_k}^*\alpha\). Since it kills
exceptional divisors and \(L_k-k\{S\}\) stays bounded,
\[
 0\leq P_k\cdot\eta_k=O(k^{-1}),\qquad
 K_{B_k}\cdot\eta_k=K_B\cdot\eta<0.
\]
For a fixed ample \(H\) on \(B\), polarize by
\(P_k+k^{-1}r_k^*H\) plus a sufficiently small rational ample
class on \(B_k\). Approximate \(\eta_k\) by positive sums of
covering-curve classes. Discard terms with nonnegative
canonical degree. Some remaining covering curve has the ratio
of polarization degree to anticanonical degree \(O(k^{-1})\).
Quantitative bend-and-break
\cite[Theorem~5]{Proj-MiyaokaMori1986} supplies rational curves
through its general points with polarization degree at most
\(2\dim B\) times that ratio. Parameterization and
uncountability give a covering family of rational curves
\(R_k\) satisfying
\begin{equation}\label{proj:ex:small-curves}
 P_k\cdot R_k=O(k^{-1}),\qquad
 H\cdot r_{k*}R_k=O(1).
\end{equation}
The general parametrized member is free in characteristic
zero. Consequently \(K_{B_k}\cdot R_k\leq-2\), and
\begin{equation}\label{proj:ex:contact-bound}
 0\leq K_{B_k/B}\cdot R_k
   =K_{B_k}\cdot R_k-K_B\cdot r_{k*}R_k=O(1).
\end{equation}
Here the upper bound follows from the fixed ample-degree
bound in \eqref{proj:ex:small-curves}.

Write \(\bar R_k=r_{k*}R_k\). Every positive contact with an
exceptional prime contributes a positive integer discrepancy
times a positive integer contact degree to
\eqref{proj:ex:contact-bound}. Thus the number of such contacts,
their discrepancies, and their contact degrees are uniformly
bounded. Contacts with the finitely many strict transforms
of divisorial Lelong components of \(S\) are bounded by their
fixed degrees against \(\bar R_k\). Subtracting these
divisorial parts from \(r_k^*S\) leaves a positive current.
Using \eqref{proj:ex:multiplier-order} and
\eqref{proj:ex:small-curves} gives
\begin{equation}\label{proj:ex:zero-gap}
 \begin{split}
 0&\leq \{S\}\cdot\bar R_k
          -\sum_E\nu_E(S)\,E\cdot R_k\\
  &\leq \{S\}\cdot\bar R_k-(F_k/k)\cdot R_k
   \leq O(k^{-1}).
 \end{split}
\end{equation}
The sum includes exceptional primes and those strict
divisorial components; terms with zero contact are irrelevant.

Bounded ample degree gives only finitely many numerical
classes of the integral cycles \(\bar R_k\). Pass to one
class. By Lemma~\ref{proj:ex:valuation-acc}, the sums in
\eqref{proj:ex:zero-gap} belong to an ACC set: there are a bounded
number of terms, their nonnegative integral multiplicities
are bounded, and all relevant discrepancies are bounded.
Finite sums of this kind preserve ACC, by successively
taking nonincreasing subsequences of their entries.
The nonnegative gaps in \eqref{proj:ex:zero-gap} must therefore
equal zero for some covering families. Otherwise a sequence
tending to zero has a strictly decreasing positive
subsequence, giving a strictly increasing sequence of the
complementary sums.

For such a family the residual positive current restricts
to degree zero on almost every \(R_k\), and hence vanishes
there. Off its removed divisors it is the original current.
The discussion following Lemma~\ref{proj:ex:web} and that lemma
itself yield a positive-relative-dimensional rational
quotient of \(B\) whose vertical directions are annihilated
by \(S\) on an open. Composing with \(h\) gives a quotient
of smaller base dimension whose vertical directions are
annihilated by \(T\), contradicting the choice of \(b\).
\end{proof}

\subsection{Excluding normalized bounded curves}

Run a canonical LMMP on \(X\) with scaling of a fixed very ample
rational divisor \(A\). For rational \(t>0\), its positive-threshold
stages give terminal \(\Q\)-factorial models \(Y_t\) on which
\(K_t+tA_t\) is nef, big, and semiample
\cite{BCHM}. There are only finitely many divisorial
contractions, since each lowers Picard number. Fix a late model
\(Y\) after the last such contraction. All subsequent maps
\(Y\dashrightarrow Y_t\) are small, and are canonical
nonpositive birational maps. If the program terminates, use
the last model repeatedly. Define
\begin{equation}\label{proj:ex:volume-scale}
 \mu_t=\vol(K_X+tA)^{1/n}.
\end{equation}
Since \(K_X\) is not big, \(\mu_t\to0\). In particular \(K_Y\)
is nonbig. The transform \(A_t\) is effective up to rational
linear equivalence; all intersections below use general
covering curves, which are not contained in a chosen such
representative.

\begin{lemma}[Exceptional current comparison]\label{proj:ex:current-comparison}
Fix a positive current \(T\) on a smooth resolution in the
pullback class of \(K_Y\). On a common smooth resolution of
\(Y\dashrightarrow Y_t\), write
\[
 p^*K_Y=p_t^*K_t+J_t,\qquad J_t\geq0,
\]
where \(J_t\) is exceptional over \(Y_t\).
Then \(\nu_E(T)\geq\operatorname{coeff}_E J_t\) for every
prime on that resolution.
\end{lemma}
\begin{proof}
Pass to a common smooth model \(V\) also dominating the
fixed resolution carrying \(T\). Choose on the fixed
resolution a sufficiently positive line bundle \(H\).
For sufficiently divisible \(k\), multiplier approximation
gives a globally generated
\(\OO(kp^*K_Y+H)\otimes\mathcal J(kT)\).
Pull this system to \(V\). A general member \(D_k\) of its
underlying effective divisor system has, at every specified
prime \(E\), multiplicity
\(\ord_E\mathcal J(kT)\); global generation after removal
of the ideal ensures no additional generic vanishing.
Here the notation \(H\) also denotes its pullback to \(V\).

Put \(D'_k=p_{t*}D_k\) and \(H_t=p_{t*}H\). The first is
effective, and both are \(\Q\)-Cartier because \(Y_t\) is
\(\Q\)-factorial. Pushing the rational linear relation gives
\[
 D'_k\sim_\Q kK_t+H_t.
\]
Thus the exceptional divisor
\(D_k-p_t^*D'_k\) is rationally equivalent to
\[
 kJ_t+H-p_t^*H_t.
\]
The two exceptional divisors are equal: their difference is
numerically trivial and exceptional, so the negativity lemma
applied with both signs makes it zero. Since
\(p_t^*D'_k\) is effective,
\[
 \operatorname{coeff}_E D_k
 \geq k\operatorname{coeff}_E J_t
       +\operatorname{coeff}_E(H-p_t^*H_t).
 \]
The last coefficient is fixed for this model and \(t\).
Combining with \eqref{proj:ex:multiplier-order}, dividing by
\(k\), and letting \(k\to\infty\) proves the assertion.
The same proof can be made on each common model, so the
comparison applies to all divisorial valuations needed
later. No uniformity in \(t\) of the fixed error is required.
\end{proof}

\begin{proposition}\label{proj:prop:bounded-curves}
There are no sequences \(t_j\downarrow0\) and covering
families of curves on \(Y_{t_j}\), with smooth projective
domains \(C_j\), for which both
\[
 g(C_j),\qquad
 \frac{(K_{t_j}+t_jA_{t_j})\cdot C_j}{\mu_{t_j}}
\]
are bounded by a fixed constant. Degrees mean degrees on
the domains, or equivalently intersection with their
pushforward cycles.
\end{proposition}
\begin{proof}
Suppose such families exist. Slice their parameters and
resolve total evaluation, the map to the fixed \(Y\), the
fixed smooth resolution carrying \(T\), and the common
models comparing \(Y\) and \(Y_t\). Denote the resulting
generically finite dominant morphism by \(q_t:U_t\to Y\).
The source is smooth near its complete general parameter
fiber and has dimension \(n\), with parameter space of
dimension \(n-1\). Rational maps from this smooth source
to each required proper target extend at codimension-one
points by the valuative criterion of properness. Their
indeterminacy loci consequently have codimension at least
two, and their images cannot dominate the parameter space.
Shrink that space to avoid these finitely many images,
and resolve and compactify preserving the remaining open.
Equivalently, nontrivial smooth-center blowups have
codimension-at-least-two centers; a codimension-one
Cartier-center blowup is the identity. Thus the complete
smooth general parameter fiber is still \(C_t\), unchanged
as a curve, and its genus is unchanged.

For a prime \(Z\) of \(U_t\) exceptional over \(Y\), the
restricted valuation of its function field is
\(r_Z\ord_E\) for a divisorial valuation \(E\) over \(Y\).
The coefficient of \(Z\) in the ramification difference is
\begin{equation}\label{proj:ex:ramification-cost}
 \operatorname{coeff}_Z(K_{U_t}-q_t^*K_Y)
       =r_Za(E;Y,0)-1.
\end{equation}
This follows by factoring through a model extracting \(E\)
and calculating the tame ramification of DVRs. It is
positive and bounded away from zero: \(Y\) is terminal,
so \(a(E;Y,0)>1\), and a fixed Cartier index of \(K_Y\)
puts discrepancies in a fixed rational lattice. The
remaining ramification coefficients are nonnegative.
On the general parameter fiber,
\[
 K_{U_t}\cdot C_t=2g(C_t)-2,\qquad
 q_t^*K_Y\cdot C_t\geq0.
 \]
The latter inequality follows from pseudoeffectivity and
the covering property. Therefore the total ramification
cost is bounded. Formula~\eqref{proj:ex:ramification-cost}
bounds the number of positive exceptional contacts,
their discrepancies, their ramification indices, and
their integral contact degrees. Moreover
\(q_t^*K_Y\cdot C_t\) lies in a fixed rational lattice
inside a bounded interval, so takes only finitely many
values.

The generic Lelong number of \(q_t^*T\) at \(Z\) is
\(r_Z\nu_E(T)\). Extract \(E\), remove its generic
divisorial part, and use the transverse-power-map
comparison for the zero-generic-number remainder.
Lemma~\ref{proj:ex:current-comparison} shows that these
coefficients dominate the coefficients of the pulled-back
canonical difference \(J_t\). Its support is exceptional
over \(Y\) as well as \(Y_t\), since the map between these
models is small. Subtracting all exceptional divisorial
parts of \(q_t^*T\) leaves a positive current. Consequently
\begin{equation}\label{proj:ex:bounded-gap}
 \begin{split}
 0&\leq q_t^*K_Y\cdot C_t
       -\sum_{Z\ {\rm exceptional}}
                r_Z\nu_E(T)\,Z\cdot C_t\\
  &\leq K_t\cdot C_t
   \leq (K_t+tA_t)\cdot C_t
   =O(\mu_t).
 \end{split}
\end{equation}
All divisor contacts used here are nonnegative, since the
curves cover the total space.

On the fixed smooth resolution \(W\to Y\), one has
\[
 a(E;W,0)=a(E;Y,0)
       -\ord_E(K_W-p^*K_Y)\leq a(E;Y,0),
\]
because the relative canonical divisor is effective.
Thus Lemma~\ref{proj:ex:valuation-acc} applies with a fixed
bound. The sums in \eqref{proj:ex:bounded-gap} form an ACC
set, their number of terms and integral multiplicities
being bounded. Pass to a fixed value of
\(q_t^*K_Y\cdot C_t\). Since \(\mu_t\to0\), the
exact-zero argument of \eqref{proj:ex:zero-gap} gives a family
for which the left gap is zero.

The residual positive current then restricts to zero on
almost every curve in that family. On the dense open
away from the removed exceptional divisors it agrees
with the original current. Its curve tangents therefore
generate, by Lemma~\ref{proj:ex:web}, an annihilated rational
fibration of positive relative dimension on \(Y\).
This contradicts Proposition~\ref{proj:prop:no-pullback}.
\end{proof}

\subsection{The signed alternative}

\begin{proposition}\label{proj:prop:signed-alternative}
On any model \(Y_t\), let a signed rational divisor
represent \(K_t\) up to \(\Q\)-linear equivalence. On a
log resolution \(p:W\to Y_t\), let \(D\) be the reduced
SNC divisor consisting of its strict support and all
exceptional divisors. Then \(K_W+D\) is big.
\end{proposition}
\begin{proof}
Put \(P=K_W+D\), and suppose it is not big. It cannot
have Iitaka dimension at least one: by
Corollary~\ref{proj:ex:supporting}, \(K_W+cP\) would be big,
and
\((1+c)P=(K_W+cP)+D\) would then be big as well.
Hence \(\kappa(W,P)\leq0\).

Run the dlt LMMP for \(P\) with ample scaling.
The divisor \(P\) has a signed rational representative
supported on the floor \(D\), and its transforms retain
that property. Every negative flip must meet the floor:
on its complement the representing rational section is
nowhere vanishing, so the adjoint has zero degree on
every complete curve there. By
Proposition~\ref{proj:prop:special-termination}, an infinite
sequence would eventually have all flips disjoint from
the floor, a contradiction. Divisorial contractions
are finite in number. Thus the program terminates at
a \(\Q\)-factorial dlt log minimal model
\((Z,D_Z)\).

The boundary is reduced, \(K_Z\) is pseudo-effective
as the birational pushforward of \(K_W\), and the nef
adjoint has a signed rational representative supported
on \(D_Z\). Its restriction to \(D_Z\) is semiample
by Proposition~\ref{proj:prop:boundary-restriction}.
Theorem~\ref{proj:thm:signed} applies, with the
pseudo-effective perturbation \(K_Z\), and gives
abundance. Since the Iitaka dimension is at most zero,
the numerical dimension is zero. Intersecting
\(K_Z+D_Z\equiv0\) with an ample
\((n-1)\)-fold product, and using pseudoeffectivity of
\(K_Z\), forces \(D_Z=0\). The signed relation then
gives \(K_Z\sim_\Q0\).

On a common resolution, the pullback of \(K_W\) is
rationally equivalent to a signed divisor exceptional
over \(Z\). It is pseudo-effective. A positive current
in that class has zero intersection with a pullback
of an ample class on \(Z\) to the power \(n-1\),
so is supported on the exceptional locus. The support
theorem makes it effective divisorial, and independence
of exceptional divisor classes by negativity says that
the original signed coefficients are these nonnegative
coefficients. Thus \(K_W\) is \(\Q\)-linearly effective,
contradicting \eqref{proj:ex:counterexample}.
\end{proof}
